\section{Discussion and Current Limits}

This draft deliberately stops short of a full interpretation layer. Still, the manuscript should state the boundaries of the current evidence clearly so that later quantitative claims remain anchored to what the simulation pipeline actually models.

The main limitations that should remain visible in the first paper version are the following:
\begin{itemize}
    \item Only a single outdoor scenario is wired into the current manuscript scaffold.
    \item The current Rabot configuration disables coverage-map exports, so the paper centers on trajectory-level \gls{sinr} and schedule metrics.
    \item The current checked scenario also disables capped exact search, so the default draft compares the baseline and the local \gls{csi}-driven heuristic unless an updated output set is committed.
    \item Citizen participation is abstracted as movable candidate positions and scheduler actions rather than being validated experimentally.
    \item The present text does not yet discuss sensitivity to seeds, scenario perturbations, or broader deployment costs.
\end{itemize}

The next analysis pass should therefore answer at least three questions: how stable the reported gains are across seeds and scenarios, how strongly the conclusions depend on the chosen \gls{sinr} objective, and how the relocation metrics should be interpreted in a real co-created deployment process.

\gilles{Later pass: split this section into limitations and future work if the venue expects a more formal discussion structure.}
